% Generated draft: requires mathematical and source review.
\documentclass[11pt]{article}
\usepackage[margin=25mm]{geometry}
\usepackage{fontspec}
\setmainfont{Latin Modern Roman}
\usepackage{amsmath,amssymb,mathtools,unicode-math}
\setmathfont{Latin Modern Math}
\usepackage{enumitem,needspace,xurl,hyperref}
\setlength{\parindent}{0pt}
\setlength{\parskip}{5pt}
\setlength{\emergencystretch}{3em}
\newcommand{\N}{\mathbb{N}}
\newcommand{\Z}{\mathbb{Z}}
\newcommand{\Q}{\mathbb{Q}}
\newcommand{\R}{\mathbb{R}}
\newcommand{\C}{\mathbb{C}}
\newcommand{\F}{\mathbb{F}}
\newcommand{\A}{\mathbb{A}}
\newcommand{\PP}{\mathbb{P}}
\newcommand{\E}{\mathbb{E}}
\newcommand{\eps}{\varepsilon}
\newcommand{\dd}{\,\mathrm{d}}
\begin{document}
\section*{UnsolvedMath ID: 10400090}
\textbf{Source problem code:} catalogue entry
\subsection{Definitions and mathematical statement}
The source asks for a construction rather than a fixed numerical identity: construct invariants of three-manifolds through a linear skein theory based on the Kauffman skein module. Once the coefficient ring \(R\), link/framing conventions and Kauffman skein relations are fixed, such a module has the general form \[ \mathcal S_K(M;R)=R\langle\mathcal L(M)\rangle/\mathcal R_K(M), \] where \(\mathcal L(M)\) is the chosen set of isotopy classes of links in \(M\) and \(\mathcal R_K(M)\) is the submodule generated by those specified local relations. This display records the module construction only; it does not select the missing Kauffman relation or parameter values. The problem seeks a linear skein-theoretic construction \(I\) from this data, with an identified target, such that for every pair of three-manifolds in the intended category, \[ M\cong N\ \Longrightarrow\ I(M)=I(N) \] (or an isomorphism of outputs, for structured targets). The relevant equivalence must preserve whatever orientations and boundary data the category requires. No scalar target, functorial field theory, gluing axiom, nontriviality criterion or effectiveness requirement is stated in the excerpt, so none is imposed here. The module's own invariance under the specified homeomorphisms does not by itself identify the additional construction meant by the source.
\subsection{Short English statement}
Develop a linear skein construction of three-manifold invariants using a Kauffman skein module. The source leaves the module conventions and the desired kind of invariant unspecified.
\subsection{Sources}
\begin{itemize}
\item Local UnsolvedMath registry, permanent ID \(10400090\), exact source code \texttt{catalogue entry}.
\item Ohtsuki, \emph{Problems on invariants of knots and 3-manifolds} (2002), Problem 4.12, PDF page 77: \url{https://msp.org/gtm/2002/04/gtm-2002-04-024s.pdf}. The registry records access on 2026-07-29 and extraction by pdf-font-structure.
\end{itemize}
\end{document}
